\subsection{A five-point Krammer quotient retaining the cubic detector}

\paragraph{D23.P6 (a concrete extension outside the low-derived tests).}
Let $k=\mathbb F_3[i]$, $i^2=-1$. In the ordered basis
$(F_{12},F_{13},F_{14},F_{23},F_{24},F_{34})$, put
\[
A=\begin{pmatrix}
-i&0&0&-1-i&-1-i&0\\
0&0&0&i&0&0\\
0&0&0&0&i&0\\
0&1&0&1-i&0&0\\
0&0&1&0&1-i&0\\
0&0&0&0&0&1
\end{pmatrix},\quad
B=\begin{pmatrix}
0&i&0&0&0&0\\
1&1-i&0&0&0&0\\
0&0&1&0&0&0\\
0&1-i&0&-i&0&-1-i\\
0&0&0&0&0&i\\
0&0&0&0&1&1-i
\end{pmatrix},
\]
\[
C=\begin{pmatrix}
1&0&0&0&0&0\\
0&0&i&0&0&0\\
0&1&1-i&0&0&0\\
0&0&0&0&i&0\\
0&0&0&1&1-i&0\\
0&0&1-i&0&1-i&-i
\end{pmatrix}.
\tag{D23.P6.1}
\]
The assignment $(\sigma_1,\sigma_2,\sigma_3)\mapsto(A,B,C)$
is a continuous finite representation of $\widehat B_4$ and
$(ABC)^4=-I$. It therefore defines an actual projective quotient
of $L=B_4/\langle z_4\rangle$, the label-five stabilizer in
$\Gamma_{0,[5]}$. Induction from that subgroup gives a finite
quotient of the full five-point group if desired.

The subspace $W=\langle F_{12},F_{13},F_{23}\rangle$ is
invariant under $A,B$. Its representation is exactly the Frobenius
contragredient of the affine detector D23.P3. More precisely, with
\[
P=\begin{pmatrix}
1-i&1-i&-1\\
-1-i&1-i&i\\
-1&-1-i&1
\end{pmatrix},
\quad
P^{-1}=\begin{pmatrix}
0&-1&i\\
-i&-1&0\\
1-i&1-i&1+i
\end{pmatrix},
\tag{D23.P6.2}
\]
and $\overline{i}=-i$, one has
\[
 A|_W\,P=P\,\overline{\eta(\sigma_1)}^{-\mathsf T},\qquad
 B|_W\,P=P\,\overline{\eta(\sigma_2)}^{-\mathsf T}.
\tag{D23.P6.3}
\]
Here affine matrices are written homogeneously in dimension three,
$\sigma_1=(sr)^{-1}$ and $\sigma_2=s\sigma_1s$.
Consequently $f(A^2,B^2)|_W$ determines the cubic sign $a_3(h)$
for every actual gentle $h=(1,f)$ by D23.P3. This is an actual
five-point quotient retaining that sign; it is not merely a
four-point forgetful product.
The same recovery works from the projective evaluation: the
homogeneous affine image $E_F$ contains no nonidentity scalar
matrix, so its affine normalization selects the unique lift on $W$.

The generated first-pair image
$P_{12}=\langle A^2,B^2\rangle$ has order $46656$, with derived
series orders
\[
46656,\quad3888,\quad486,\quad243,\quad3,\quad1.
\tag{D23.P6.4}
\]
In particular $P_{12}''$ is not nilpotent. Its quotient on $W$
is D23.P3's $E_F$, whose second derived group is
$k^2\rtimes\langle-I\rangle$, a nonnilpotent group of order $162$.
Thus the nilpotent-second-derived exclusion criterion of supplied
FP.36 does not dispose of this representation.

\begin{proof}
The matrices are the Krammer matrices at $q=i,t=-i$ in the
corrected six-case formula of Bigelow, \emph{The Lawrence--Krammer
representation}, arXiv:math/0204057v1, section4.2, printed page13.
No assertion identifying this specialized lattice with the
homological Lawrence lattice is used. Alternatively, the three
braid/distant relations and $(ABC)^4=-I$ follow by direct
multiplication of (D23.P6.1); the exact script
\path{agents/pentagon/krammer_f9_probe.py} verifies them.
All matrices are invertible, as follows either from the formula
or from their explicitly invertible diagonal blocks after the
corresponding basis permutation. Hence they define the asserted
finite representation, which extends continuously to the profinite
completion. Killing the scalar full twist is precisely the
projective passage needed for $L$; no spherical relation is
assumed for an unsupported direct $B_5$ representation.

Inspection of the first, second and fourth columns shows the
invariance of $W$. Its two matrices are
\[
K_1=\begin{pmatrix}-i&0&-1-i\\0&0&i\\0&1&1-i\end{pmatrix},
\qquad
K_2=\begin{pmatrix}0&i&0\\1&1-i&0\\0&1-i&-i\end{pmatrix}.
\]
Direct multiplication verifies both $PP^{-1}=I$ and
(D23.P6.3). In encoded $\mathbb F_9$ arithmetic these are
checked independently in
\path{agents/pentagon/krammer_affine_conjugacy.py}.
Since Frobenius followed by contragredience is an invertible
operation on matrix groups, these identities identify the entire
marked $B_3$ representation on $W$, rather than merely its order
or composition factors. Thus its pure image is exactly the
Frobenius contragredient of $E_F$, and its value on $f$ recovers
$\eta(f)$ by
\[
 \eta(f)=\overline{\bigl(P^{-1}f(A^2,B^2)|_W P\bigr)^{-\mathsf T}}.
\]
The cubic-sign formula now follows from D23.P3.

For completeness, the finite-order statements are certified in two
ways. The script \path{agents/pentagon/krammer_f9_probe.py}
uses a faithful permutation action on a union of vector orbits
containing a basis, and computes its stabilizer chain and derived
series. The independent matrix enumeration
\path{agents/pentagon/krammer_first_pair.py} enumerates exactly
$46656$ distinct matrices closed under the two generators and
exactly $3888$ matrices generated by the normal closure of their
commutator. The latter commutator has $486$ conjugates and three
of them generate the derived subgroup. The associated finite data
are retained in the JSON/NPZ certificates. These checks agree.

Finally, $E_F'=k^2\rtimes Q_8$ and
$E_F''=k^2\rtimes C_2$, since $-I$ acts as $-1$ on $k^2$.
The latter group is nonnilpotent: its Sylow two subgroup does not
commute with its normal Sylow three subgroup. A quotient of a
nilpotent finite group is nilpotent. The surjection
$P_{12}''\twoheadrightarrow E_F''$ therefore proves the final
assertion independently of the complete derived-series counts.
\end{proof}

\paragraph{Current scope.}
D23.P6 constructs and checks a new finite five-point test. It does
not prove the pentagon in this quotient for every gentle element,
and it does not prove $\Delta|_{\widehat{GT}}=0$ or exhibit a
full-GT counterexample. The exact first-pair evaluations are further
reduced to six matrices in D23.F13. The correlations with the other
two pentagon pair types remain the substantive calculation.
